

%% ===================================================================
\section{관측성(Observability) 이론}
\label{sec:lit-observability}
%% ===================================================================

\chg{정적 RTO의 한계를 극복하려면 복구 환경의 상태 변화를 실시간으로 파악할 수 있어야 한다. 이러한 상태 파악 능력을 이론적으로 뒷받침하는 개념이 관측성(observability)이며, 이 절에서는 제어이론에서 출발한 관측성 개념과 그 현대적 확장을 검토한다.}

%% -------------------------------------------------------------------
\subsection{제어 이론에서의 관측성 (Kalman, 1960)}
\label{subsec:kalman-observability}
%% -------------------------------------------------------------------

관측성(observability)의 개념은 \citet{kalman1960control}이 제어 이론에서 체계적으로 형식화하였다.
Kalman은 선형 시불변(Linear Time-Invariant, LTI) 시스템
\begin{equation}
  \dot{x}(t) = Ax(t) + Bu(t), \quad y(t) = Cx(t)
  \label{eq:lti-system}
\end{equation}
에서 시스템의 출력 $y(t)$를 관찰함으로써 초기 상태 $x(0)$를 유일하게 결정할 수 있는지를
판별하는 기준으로 관측성 행렬(observability matrix)의 계수(rank) 조건을 제시하였다.
관측성 행렬
\begin{equation}
  \mathcal{O} = \begin{bmatrix} C \\ CA \\ CA^2 \\ \vdots \\ CA^{n-1} \end{bmatrix}
  \label{eq:obs-matrix}
\end{equation}
이 완전 계수(full rank)일 때 시스템이 완전 관측 가능(completely observable)하다.

이 개념은 시스템의 내부 상태를 외부 출력만으로 추론할 수 있는 역량을 수학적으로
정의한 것으로, 단순한 모니터링(monitoring)과는 본질적으로 구별된다~\citep{beyer2016sre}.
모니터링이 사전에 정의된 지표의 임계값 초과 여부를 감시하는 반응적(reactive) 방식인 반면,
관측성은 시스템 출력의 충분성(sufficiency)에 관한 구조적 속성이다.

\chg{Kalman 이후 제어 이론에서 관측성은 상태 추정(state estimation)의 토대로 발전하였다. \citet{luenberger1966observers}는 관측성 조건이 충족되는 선형 시스템에서 직접 측정되지 않는 내부 상태를, 가용한 출력과 입력만으로 구동되는 보조 동역학계인 상태 관측기(state observer)를 통해 재구성할 수 있음을 보였다. 이로써 관측성은 가부 판별을 넘어 관측 가능한 시스템에서는 내부 상태를 실제로 복원할 수 있다는 구성적 의미를 획득하였고, 이후 상태 추정과 피드백 제어의 핵심 원리로 자리 잡았다. 본 연구가 관측성을 복구 환경의 내부 상태에 대한 가시성으로 해석하는 것은, 바로 이 출력으로부터 내부 상태를 복원하는 역량이라는 제어 이론적 함의를 BCM 영역으로 확장한 것이다.}


%% -------------------------------------------------------------------
\subsection{BCM 영역으로의 확장}
\label{subsec:observability-bcm}
%% -------------------------------------------------------------------

%% --- BCM 맥락의 정규화 ---
\subsubsection{BCM 맥락에서의 관측성 정규화}
\label{subsubsec:obs-normalization}

본 연구에서는 Kalman의 제어 이론적 관측성 개념을 BCM 영역으로 확장하여 적용한다.
\chg{Kalman의 원래 정식화는 관측성 행렬의 계수 조건을 통해 시스템이 ``완전 관측
가능한가''를 가능과 불가능의 이분법으로 판별한다. 그러나 \chg{BCM과 재해복구} 실무에서 조직의 관측성은 이분법으로 주어지지 않는다.
\chg{현대 IT 운영에서 관측성은 로그(log), 메트릭(metric), 트레이스(trace)라는 세 가지
텔레메트리(telemetry)를 통해 확보되며~\citep{majors2022observability}, 조직마다 센서와
모니터링 도구의 구비 정도, 그리고 이 세 텔레메트리의 수집과 보존 수준이 달라 그 확보
정도가 연속적으로 변한다.} 그 결과 조직의 관측성은 0과 1 사이에서 연속적으로 분포한다. 본 연구는 바로 이 부분을 확장하여 적용하는데, Kalman의
이진 판별 대신 관측성의 \emph{정도}를 $O_{\text{raw}} \in [0, 1]$의 연속 점수로 정량화한다.
이렇게 연속화하는 이유는, 복구 시간의 단축 효과가 관측성의 유무가 아니라 그
수준에 연속적으로 의존한다고 보기 때문이며, 연속 점수라야 시그모이드 바운딩을 통한 복구 시간
단축량의 정량적 모형화가 가능해지기 때문이다.}
이때 $O_{\text{raw}} = 0$은 시스템 상태를 전혀 파악할 수 없는 상태를,
$O_{\text{raw}} = 1$은 완전히 관측 가능한 상태를 의미한다.


%% --- 관측성-복구 능력 구조적 연결 ---
\subsubsection{관측성과 복구 능력의 구조적 연결}
\label{subsubsec:obs-recovery-link}

\citet{avizienis2004basic}와 \citet{hollnagel2006resilience}이 각각
의존성과 복원공학 문헌에서 논의한 바와 같이, 관측성과 복구 능력 사이에는
구조적 연결이 존재한다.
시스템 상태를 보다 정확히 파악할 수 있을수록 장애의 근본 원인(root cause)을
신속히 진단할 수 있고, 이에 따라 복구에 소요되는 시간이 단축된다.
\citet{avizienis2004basic}는 가용성(availability), 신뢰성(reliability),
안전성(safety), 무결성(integrity)을 포괄하는 의존성(dependability)의 분류 체계에서
장애 진단 능력이 시스템의 가용성에 직접적으로 영향을 미침을 논의하였다.
\citet{hollnagel2006resilience}는 복원공학(Resilience Engineering) 관점에서 시스템의
예측, 감시, 대응, 학습 능력을 강조하며, 이 중 감시(monitoring) 능력이 본질적으로
관측성에 해당한다고 볼 수 있다.
\chg{본 연구는 이 두 이론을 관측성과 복구 시간 사이의 단조 감소 관계로 반영하여 적용한다.
\citet{avizienis2004basic}에서는 장애 진단 능력이 가용성에 기여한다는 인과를,
\citet{hollnagel2006resilience}에서는 감시 능력이 복원력의 핵심이라는 명제를 각각
가져온다. 본 모델은 이 ``관측성에서 진단으로, 진단에서 복구 시간 단축으로''의 인과를
관측성 채널의 하향 조정량 $E_{O,\text{bnd}} \leq 0$으로 구현한다. 즉 관측성이 높을수록
$\DRTO$를 기준 RTO $R_0$ 아래로 끌어내리는 음의 효과로 정량화하는 것이다. 두 이론을
이렇게 적용한 이유는, 이론이 제시한 인과 관계를 모델에서 관측성과 복구 시간 사이의
단조 감소 관계로 직접 반영함으로써, 관측성 투자에 따른 복구 시간 단축을 정량적으로
예측할 수 있게 하기 위함이다.}


%% --- SRE와 세 가지 핵심 요소 ---
\subsubsection{\chg{현대 IT 운영의 관측성 SRE와 세 가지 핵심 요소(three pillars)}}
\label{subsubsec:obs-sre-pillars}

현대 IT 운영에서의 관측성은 Kalman의 제어 이론적 개념을 분산 시스템
환경으로 재해석한 것으로 볼 수 있으며, \citet{beyer2016sre}와
\citet{majors2022observability}이 이를 SRE 실무 체계로 정착시켰다.
\citet{beyer2016sre}는 Google의 SRE(Site Reliability Engineering)
실무를 통해 모니터링과 관측성의 차이를 구분하였고,
\citet{majors2022observability}는 관측성의 세 가지 핵심 구성 요소(three pillars), \chg{곧
메트릭(metrics)과 로그(logs), 트레이스(traces)를} 체계화하였다.
이 세 구성 요소의 통합적 활용을 통해 조직은 예상치 못한 장애 상황에서도
시스템의 내부 상태를 효과적으로 파악할 수 있다.
\p{[}그림~\ref{fig:obs-three-pillars}\p{]}는 이러한 세 가지 핵심 요소(three pillars)의 구성 관계를 도식화한 것이다.


%% --- 운영화 매핑 ---
\subsubsection{\chg{데이터 유형에서} \chg{도구와 목적} \chg{차원으로의 운영화 매핑}}
\label{subsubsec:obs-operationalization}

이러한 세 가지 핵심 요소(three pillars) 분류는 관측성을 구성하는 \textbf{데이터 유형(data primitives)} 차원의
정의이다. 본 연구는 \chg{BCM과 재해복구} 맥락의 운영화 단계에서 관측성을 \chg{도구와 목적}별 통합 시스템
차원, \chg{곧 APM(Application Performance Monitoring, 서비스 신뢰성 영역)과 SIEM(Security
Information and Event Management, 보안 이벤트 대응 영역)의 두 축으로} 재구성하여 정량화한다.
APM과 SIEM은 각각 \chg{메트릭, 로그, 트레이스}를 통합 활용하되 측정 목적과 주된 데이터 출처가
상이하므로 분리 산출 후 가중 합성($O_{\text{raw}} = \alpha \cdot O_{\text{APM}}
+ (1-\alpha) \cdot O_{\text{SIEM}}$)으로 통합하며, 이론적 분류와 운영화 모델의 매핑 관계는
본문 \secref{app:operationalization}(제5장 실무적 기여)의 조작화 예시에서 참고용으로 다룬다.

\begin{figure}[htbp]
\centering
\begin{tikzpicture}[
  every node/.style={text=black},
  pillar/.style={draw, thick, rounded corners=4pt, fill=blue!10, minimum width=40mm,
    minimum height=22mm, font=\Large\bfseries, align=center, text=black},
  obs/.style={draw, very thick, rounded corners=4pt, fill=green!15, minimum width=64mm,
    minimum height=18mm, font=\Large\bfseries, align=center, text=black},
  arr/.style={->, >=Stealth, very thick, black}
]
  \node[pillar] (m) at (0, 0)   {Metrics\\\normalsize\mdseries(메트릭)};
  \node[pillar] (l) at (5.5, 0) {Logs\\\normalsize\mdseries(로그)};
  \node[pillar] (t) at (11, 0)  {Traces\\\normalsize\mdseries(트레이스)};
  \node[obs]    (o) at (5.5, -3.4)
    {Observability\\\normalsize\mdseries(관측성)};
  \draw[arr] (m.south) -- (o.north);
  \draw[arr] (l.south) -- (o.north);
  \draw[arr] (t.south) -- (o.north);
\end{tikzpicture}
\caption{관측성의 세 가지 핵심 구성 요소(three pillars)}
\label{fig:obs-three-pillars}
\end{figure}

관측성 수준이 높은 조직은 장애 발생 시 원인 파악에 소요되는 시간(Mean Time to Detect,
MTTD)과 복구에 소요되는 시간(Mean Time to Recovery, MTTR)을 모두 단축할 수 있다~\citep{forsgren2018accelerate, kim2016devops}.
이는 관측성이 단순한 기술적 역량을 넘어 조직의 복구 목표 달성에 직접적으로
기여하는 운영 변수임을 시사한다.
본 연구에서는 이러한 이론적 근거를 바탕으로 관측성 점수 $O_{\text{raw}} \in [0,1]$을
$\DRTO$ 모델의 핵심 입력 변수로 도입한다.


%% -------------------------------------------------------------------
\subsection{모니터링과 관측성의 개념적 구분}
\label{subsec:monitoring-vs-observability}
%% -------------------------------------------------------------------

실무에서 모니터링(monitoring)과 관측성(observability)은 혼용되는 경우가 많으나,
BCM 맥락에서 양자는 분석 범위와 목적이 구별된다. 모니터링은 사전에 정의된
지표를 수집하여 임계값 초과 여부를 감시하는 활동인 반면, 관측성은 시스템 외부
출력만으로 내부 상태를 추정할 수 있는 시스템 고유의 속성이다. 즉 모니터링이
관측성을 구현하기 위한 수단(도구와 프로세스)이라면, 관측성은 그 수단이 확보하는
결과(시스템의 파악 가능 정도)에 해당한다.
\p{[}표~\ref{tab:monitoring-vs-observability}\p{]}은 두 개념의 핵심 차이를 비교한 것이다.

\begin{table}[htbp]
\centering
\caption{모니터링과 관측성의 비교}
\label{tab:monitoring-vs-observability}
\small
\begin{tabularx}{\textwidth}{l X X}
\toprule
\textbf{비교 기준} & \textbf{모니터링 (Monitoring)} & \textbf{관측성 (Observability)} \\
\midrule
정의 &
  사전에 정의된 지표(메트릭, 로그, 알림 규칙)를 수집·추적하는 활동 &
  시스템 외부 출력만으로 내부 상태를 추정할 수 있는 시스템 고유 속성 \\
\addlinespace
관점 &
  ``무엇을 측정할 것인가'' (도구·프로세스 중심) &
  ``시스템이 얼마나 파악 가능한가'' (시스템 능력 중심) \\
\addlinespace
범위 &
  사전 정의된 이상 징후 탐지 (known-unknowns) &
  예상하지 못한 장애의 근본 원인 추적 포함 (unknown-unknowns) \\
\addlinespace
BCM 역할 &
  장애 감지 $\to$ 알림 발생 $\to$ 대응 절차 개시 &
  복구 과정에서 시스템 상태 가시성 확보 $\to$ 복구 시간 단축 \\
\addlinespace
$\DRTO$ 대응 &
  모니터링 도구의 존재 자체는 $O_{\mathrm{raw}}$의 필요조건 &
  $O_{\mathrm{raw}} \in [0,1]$로 정량화, 관측성 채널($E_{O,\mathrm{bnd}}$)에 직접 반영 \\
\bottomrule
\end{tabularx}
\end{table}

$\DRTO$ 모델에서는 모니터링 인프라의 수준이 $O_{\mathrm{raw}}$ 값으로 변환되어
관측성 채널에 반영되므로, 모니터링 투자가 곧 $\DRTO$ 감소로 이어지는 정량적
경로를 제공한다. 이러한 관측성의 차이는 실무 환경에서도 뚜렷하게 나타난다.
관측성이 높은($O_{\mathrm{raw}} \approx 1$) 환경의 전형적 예는
실시간 서버 모니터링 시스템(Prometheus, Grafana)이 구축된 IT 인프라로,
장애 발생 시 원인을 수 분 내에 특정하여 복구 시간을 단축하며,
금융권의 실시간 트랜잭션 모니터링이나 SRE 조직의 가시성 대시보드가 여기에 해당한다.
반대로 관측성이 낮은($O_{\mathrm{raw}} \approx 0$) 환경에서는
센서가 미비한 노후 제조 설비나 수동 점검에 의존하는 중소기업 IT 환경처럼
장비 고장의 근본 원인 파악에 수 시간이 소요되어 복구가 지연된다.

\chg{이처럼 관측성은 복구 환경의 가시성을 좌우하여 복구 시간에 직접 영향을 미치지만, 복구 시간에 내재한 변동성 자체를 설명하지는 못한다. 다음 절에서는 이 변동성을 정량화하는 불확실성 모델링을 검토한다.}
